\label{ch:nsi-theory}

This chapter develops the physics of propagation Non-Standard Interactions
(NSI) from the coherent-forward-scattering matter potential up to the
fermion-composition dependence of the effective coupling matrix. It is
deliberately implementation-agnostic: \cref{ch:nsi-implementation} maps every
equation derived here onto the concrete \texttt{tpeanuts} code that evaluates
it.

\section{The Standard Model matter potential}

\begin{theorybox}
A neutrino propagating through ordinary matter undergoes coherent forward
scattering off the electrons, protons, and neutrons of the medium. For
$\nu_e$, an additional charged-current (CC) diagram (via $W$ exchange with
electrons) is possible on top of the neutral-current (NC) diagram common to
all three active flavours, giving the Wolfenstein matter
potential~\parencite{Wolfenstein1978,MikheyevSmirnov1985}
\begin{equation}
  V_{CC} = \pm\sqrt{2}\,G_F\,n_e,
\end{equation}
with the sign fixed by the neutrino/antineutrino convention ($+$ for
neutrinos, $-$ for antineutrinos). In the flavour basis $(\nu_e,\nu_\mu,\nu_\tau)$
this contributes $\diag(V_{CC},0,0)$ to the propagation Hamiltonian: the NC
piece is common to all three flavours and is therefore an unobservable global
phase in the pure Standard Model, 3-flavour case (it becomes observable only
once a fourth, NC-blind sterile state is introduced -- a separate BSM
extension, not covered by this document).
\end{theorybox}

\section{Non-Standard Interactions: the effective $\epsilon$ matrix}

\begin{theorybox}
NSI generalise the coherent-forward-scattering picture by allowing new,
model-independent four-fermion operators of the form
\begin{equation}
  \mathcal{L}_{\rm NSI} = -2\sqrt2\,G_F\,\epsilon_{\alpha\beta}^{f,P}
  \left(\bar\nu_\alpha\gamma^\mu P_L\nu_\beta\right)
  \left(\bar f\gamma_\mu P\, f\right),
  \qquad P\in\{L,R\},
\end{equation}
where $f$ is a matter fermion (electron, up quark, or down quark) and
$\epsilon_{\alpha\beta}^{f,P}$ is a dimensionless coupling, generally complex
for $\alpha\neq\beta$~\parencite{Grossman1995}. When a neutrino of flavour
$\alpha$ coherently forward-scatters off a medium containing many such
fermions, the sum over $f$ and $P$ collapses into a single effective coupling
per flavour pair, and the propagation Hamiltonian in the flavour basis
becomes
\begin{equation}
  \keyresult{H_{\rm mat}^{\rm NSI} = V_{CC}\left(\begin{pmatrix}1&0&0\\0&0&0\\0&0&0\end{pmatrix} + \epsilon\right)},
\end{equation}
with $\epsilon$ a dimensionless $3\times3$ \emph{Hermitian} matrix,
\begin{equation}
  \epsilon = \begin{pmatrix}
    \epsilon_{ee} & \epsilon_{e\mu} & \epsilon_{e\tau} \\
    \epsilon_{e\mu}^* & \epsilon_{\mu\mu} & \epsilon_{\mu\tau} \\
    \epsilon_{e\tau}^* & \epsilon_{\mu\tau}^* & \epsilon_{\tau\tau}
  \end{pmatrix}.
  \label{eq:epsilon-matrix}
\end{equation}
Diagonal entries are real (Hermiticity); off-diagonal entries are generally
complex, each carrying its own CP-violating phase independent of the
standard PMNS $\delta$. The Standard Model limit is $\epsilon=0$.

Hermiticity of $\epsilon$ is not a modelling choice: it is required for
$H_{\rm mat}^{\rm NSI}$ itself to be Hermitian, and hence for the resulting
time evolution to be unitary (probability-conserving). Only the traceless
part of $\epsilon$ is observable in oscillation probabilities -- adding
$cI_3$ to $\epsilon$ for any real $c$ shifts $H_{\rm mat}^{\rm NSI}$ by
$c\,V_{CC}\,I_3$, and because $I_3$ commutes with every other term in the
Hamiltonian at every point along the trajectory, this generalised-time-ordered
exponential factorises into an overall, position-dependent but
flavour-independent phase, exactly as for the plain NC term above. Fixing
$\epsilon_{\mu\mu}=0$ is therefore a valid convention for the diagonal
sector, not a loss of generality.
\end{theorybox}

\begin{notebox}
\textbf{Antineutrinos.} The matter background is not its own antimatter, so
the effective potential seen by an antineutrino uses the complex-conjugated
matrix, $\epsilon\to\epsilon^{*}$, on top of the usual $V_{CC}\to-V_{CC}$ sign
flip -- exactly mirroring the $U\to U^{*}$ convention already used for the
PMNS mixing matrix itself. Off-diagonal CP-violating phases therefore flip
sign for antineutrinos, as physically required.
\end{notebox}

\section{Fermion decomposition and the electron/proton-versus-neutron split}
\label{sec:theory-composition}

\begin{theorybox}
The single matrix $\epsilon$ of \cref{eq:epsilon-matrix} is not itself a
fundamental object: it is a matter-composition-weighted combination of three
independent, per-fermion couplings $\epsilon^{e}$, $\epsilon^{u}$,
$\epsilon^{d}$~\parencite{OhlssonNSIReview2013,FarzanTortolaNSIReview2018,KikuchiMinakataUchinami2009}.
Writing $n_e(x)$, $n_p(x)$, $n_n(x)$ for the local electron, proton, and
neutron number densities, the effective coupling actually entering the
Hamiltonian at position $x$ is
\begin{equation}
  \epsilon_{\alpha\beta}(x) = \epsilon_{\alpha\beta}^{e}
  + \frac{n_p(x)}{n_e(x)}\,\epsilon_{\alpha\beta}^{p}
  + \frac{n_n(x)}{n_e(x)}\,\epsilon_{\alpha\beta}^{n},
  \qquad
  \epsilon^{p}\equiv2\epsilon^{u}+\epsilon^{d},
  \quad
  \epsilon^{n}\equiv\epsilon^{u}+2\epsilon^{d},
\end{equation}
where $\epsilon^p$/$\epsilon^n$ are themselves fixed linear combinations of
the up- and down-quark couplings (a proton is $uud$, a neutron $udd$).
Ordinary matter is electrically neutral, $n_p(x)\equiv n_e(x)$ at every
point, so the three-fermion sum collapses to exactly \emph{two} effective
pieces rather than three independent physical ones:
\begin{equation}
  \keyresult{\epsilon_{\alpha\beta}(x) = \epsilon_{\alpha\beta}^{(e+p)}
  + \frac{n_n(x)}{n_e(x)}\,\epsilon_{\alpha\beta}^{n}},
  \qquad
  \epsilon^{(e+p)} \equiv \epsilon^{e} + 2\epsilon^{u} + \epsilon^{d}.
  \label{eq:epsilon-composition}
\end{equation}
The label ``electron/proton-normalized'' for $\epsilon^{(e+p)}$ is therefore a
statement about which density it multiplies ($n_e(x)$, via $n_p(x)\equiv
n_e(x)$), not a claim that it is itself a single fundamental coupling -- it
is the fixed combination $\epsilon^{e}+2\epsilon^{u}+\epsilon^{d}$ that
happens to be proportional to $n_e(x)$ in any electrically neutral medium.
\end{theorybox}

\begin{notebox}
\textbf{Published bounds are single-composition fits, not fundamental
couplings.} Every published NSI bound or best-fit value -- including every
preset in \cref{ch:nsi-observables} -- is a measurement of the
\emph{effective}, single-medium-composition $\epsilon$ of
\cref{eq:epsilon-matrix}, typically for Earth-like matter
($n_n(x)/n_e(x)\approx1$), not of $\epsilon^{e}$, $\epsilon^{u}$,
$\epsilon^{d}$ individually. Setting $\epsilon^{n}=0$ (i.e.\ using only
\cref{eq:epsilon-matrix}, without \cref{eq:epsilon-composition}'s second
term) is therefore a \emph{modelling hypothesis} carried by whichever
preset or analysis quotes the bound, not a physical fact derivable from the
quoted number: two different experiments with different target composition
can each report an ``$\epsilon_{ee}$'' consistent with a common
$(\epsilon^e,\epsilon^u,\epsilon^d)$ but numerically different
$\epsilon^{(e+p)}$, and neither measurement alone fixes $\epsilon^n$. This
composition dependence is also the reason the LMA-Dark degenerate solution
(\cref{sec:theory-lma-dark}) is sometimes called into question when
different target compositions are combined in a single global
fit~\parencite{ColomaEtAl2017DarkSide}.
\end{notebox}

\begin{theorybox}
\textbf{Magnitude of the composition dependence.} $n_n(x)/n_e(x)$ is not a
universal number -- it is a genuine function of position within a given
medium, and it varies sharply between media relevant to
\texttt{tpeanuts}:
\begin{itemize}
  \item \textbf{The Sun.} The B16 standard solar
    model~\parencite{Vinyoles2017B16} gives a hydrogen-dominated envelope
    where $n_n(r)/n_e(r)\to0$ (a proton has no neutrons to contribute), rising
    smoothly to $n_n(r)/n_e(r)\approx1$ in the helium/metal-enriched core,
    where a non-negligible fraction of nucleons are bound in $^4$He and
    heavier nuclei with roughly equal proton and neutron content.
  \item \textbf{The Earth.} The Preliminary Reference Earth Model
    (PREM)~\parencite{DziewonskiAnderson1981} gives $n_n(x)/n_e(x)$ in the
    range $\sim1.0$--$1.15$ across the mantle/core boundary, reflecting the
    modest neutron excess of common terrestrial nuclei (silicates in the
    mantle, iron in the core) relative to a strictly isoscalar medium.
\end{itemize}
A single constant $\epsilon$ (i.e.\ $\epsilon^n=0$) is therefore an excellent
approximation for propagation confined to one roughly isoscalar medium --
which is why it was the only term implemented for most of this project's
history -- but is not exact across media of very different composition, and
is a comparatively crude approximation for propagation through the Sun's
envelope specifically, where $n_n(r)/n_e(r)$ swings from $\approx0$ to
$\approx1$ over the trajectory.
\end{theorybox}

\section{The 3+1 sterile extension and its neutral-current term}

\begin{theorybox}
NSI can be combined freely with the independent 3+1 sterile-neutrino
extension (a fourth, gauge-singlet mass eigenstate with its own mixing
angles $\theta_{14},\theta_{24},\theta_{34}$ and mass splitting
$\Delta m^2_{41}$). Because the sterile state is an
$\mathrm{SU}(2)_L\times\mathrm{U}(1)_Y$ singlet, it does not receive the
$V_{CC}$ term, but it also does not receive the neutral-current potential
$V_{NC}=\mp\frac{1}{\sqrt2}G_F\,n_n$ that the three active flavours share
equally. In the pure 3-flavour case that shared $V_{NC}$ is an unobservable
common phase (\S1.1); once a sterile state is present, subtracting it as an
overall $V_{NC}\,I_4$ phase leaves a genuine relative potential on the
sterile diagonal entry,
\begin{equation}
  \diag(V_{CC}+V_{NC},\,V_{NC},\,V_{NC},\,0) - V_{NC}\,I_4
  = \diag(V_{CC},\,0,\,0,\,-V_{NC}).
\end{equation}
This sterile neutral-current (NC) term is conceptually and mathematically
independent of the NSI composition term of
\cref{eq:epsilon-composition}: both need the local neutron density
$n_n(x)$, but they are additive, independently switchable contributions
(\cref{ch:nsi-implementation}), and the sterile NC term applies only for a
4-flavour scenario, while the NSI composition term applies for any flavour
count.
\end{theorybox}

\section{The LMA-Dark degeneracy}
\label{sec:theory-lma-dark}

\begin{theorybox}
The standard LMA solar solution ($\theta_{12}\approx33.4^\circ$) admits a
discrete degeneracy in global oscillation data: replacing
$\theta_{12}\to\pi/2-\theta_{12}$ ($\approx56.6^\circ$) together with
\begin{equation}
  \epsilon_{ee} \to -(2+\epsilon_{ee}) \approx -2.0
\end{equation}
reproduces an equivalent fit to solar and reactor data (LMA-Dark;
\parencite{Esteban2018LMADark}). The large negative $\epsilon_{ee}$
effectively flips the sign of the $\nu_e$ matter potential, so that the
octant-flipped mixing angle produces the same resonance structure as the
standard solution with the standard sign. This is a striking illustration
that a sufficiently large NSI coupling can hide an entire discrete ambiguity
in the standard oscillation parameters, and is the physics case validated
numerically in \cref{ch:nsi-results}.
\end{theorybox}

\begin{refbox}
Original NSI proposal: \cite{Grossman1995}. Model-independent one-parameter
bounds: \cite{BiggioBlennowFM2009}. Per-fermion decomposition and matter
composition dependence: \cite{OhlssonNSIReview2013,FarzanTortolaNSIReview2018,KikuchiMinakataUchinami2009}.
LMA-Dark degeneracy and its composition sensitivity:
\cite{Esteban2018LMADark,ColomaEtAl2017DarkSide}. Standard matter potential:
\cite{Wolfenstein1978,MikheyevSmirnov1985}. Solar and Earth composition
profiles: \cite{Vinyoles2017B16,DziewonskiAnderson1981}.
\end{refbox}
